% ============================================================================
% EKLER (APPENDIX)
% ============================================================================

\appendix

\chapter{EKLER}
\label{chap:ekler}

\section{Ornek Kod Parcalari}
\label{sec:ornek-kodlar}

\subsection{Kimlik Dogrulama Modulu}
\label{subsec:auth-kod}

\begin{lstlisting}[language=JavaScript, caption={auth.js - Kullanici Giris Islemi}]
const express = require('express');
const bcrypt = require('bcryptjs');
const db = require('../database/db');

const router = express.Router();

// POST /api/auth/login - Giris yap
router.post('/login', (req, res) => {
    try {
        const { username, password } = req.body;

        // Alanlar bos mu kontrol et
        if (!username || !password) {
            return res.status(400).json({ 
                error: 'Kullanici adi ve sifre gerekli' 
            });
        }

        // Kullaniciyi veritabanindan bul
        const user = db.prepare(`
            SELECT u.*, f.name as faculty_name
            FROM users u
            LEFT JOIN faculties f ON u.faculty_id = f.id
            WHERE u.username = ? AND u.is_active = 1
        `).get(username);

        if (!user) {
            return res.status(401).json({ 
                error: 'Kullanici bulunamadi' 
            });
        }

        // Sifre dogrulamasi
        const validPassword = bcrypt.compareSync(
            password, 
            user.password_hash
        );
        
        if (!validPassword) {
            return res.status(401).json({ 
                error: 'Sifre hatali' 
            });
        }

        // Session olustur
        req.session.user = {
            id: user.id,
            username: user.username,
            role: user.role,
            faculty_id: user.faculty_id
        };

        res.json({
            message: 'Giris basarili',
            user: req.session.user
        });

    } catch (error) {
        res.status(500).json({ 
            error: 'Giris yapilamadi' 
        });
    }
});

module.exports = router;
\end{lstlisting}

\subsection{Middleware - Yetkilendirme Kontrolu}
\label{subsec:middleware-kod}

\begin{lstlisting}[language=JavaScript, caption={auth.js - Yetkilendirme Middleware}]
// Oturum kontrolu middleware
const requireAuth = (req, res, next) => {
    if (!req.session.user) {
        return res.status(401).json({ 
            error: 'Oturum acilmamis' 
        });
    }
    next();
};

// Super Admin yetkisi kontrolu
const requireSuperAdmin = (req, res, next) => {
    if (req.session.user.role !== 'super_admin') {
        return res.status(403).json({ 
            error: 'Super Admin yetkisi gerekli' 
        });
    }
    next();
};

// Fakulte Admin veya ustu yetki kontrolu
const requireFacultyAdmin = (req, res, next) => {
    const allowedRoles = ['super_admin', 'faculty_admin'];
    
    if (!allowedRoles.includes(req.session.user.role)) {
        return res.status(403).json({ 
            error: 'Yetkiniz yok' 
        });
    }
    next();
};

module.exports = { 
    requireAuth, 
    requireSuperAdmin, 
    requireFacultyAdmin 
};
\end{lstlisting}

\section{Veritabani Semasi}
\label{sec:veritabani-semasi}

\begin{lstlisting}[language=SQL, caption={schema.sql - Veritabani Semasi}]
-- FAKULTELER
CREATE TABLE IF NOT EXISTS faculties (
    id INTEGER PRIMARY KEY AUTOINCREMENT,
    name TEXT NOT NULL,
    code TEXT UNIQUE NOT NULL,
    color TEXT DEFAULT '#1e40af',
    is_active INTEGER DEFAULT 1,
    created_at DATETIME DEFAULT CURRENT_TIMESTAMP
);

-- KULLANICILAR
CREATE TABLE IF NOT EXISTS users (
    id INTEGER PRIMARY KEY AUTOINCREMENT,
    faculty_id INTEGER REFERENCES faculties(id),
    username TEXT UNIQUE NOT NULL,
    password_hash TEXT NOT NULL,
    full_name TEXT NOT NULL,
    role TEXT CHECK(role IN (
        'super_admin', 'faculty_admin', 
        'supervisor', 'staff'
    )) NOT NULL,
    is_active INTEGER DEFAULT 1,
    created_at DATETIME DEFAULT CURRENT_TIMESTAMP
);

-- GOREVLER
CREATE TABLE IF NOT EXISTS tasks (
    id INTEGER PRIMARY KEY AUTOINCREMENT,
    location_id INTEGER REFERENCES locations(id),
    assigned_to INTEGER REFERENCES users(id),
    status TEXT CHECK(status IN (
        'pending', 'in_progress', 'completed', 
        'approved', 'rejected'
    )) DEFAULT 'pending',
    due_date DATE,
    created_at DATETIME DEFAULT CURRENT_TIMESTAMP
);
\end{lstlisting}

\section{Ornek Kullanici Hesaplari}
\label{subsec:ornek-kullanicilar}

\begin{table}[H]
\centering
\caption{Test Kullanici Hesaplari}
\begin{tabular}{|l|l|l|l|}
\hline
\textbf{Kullanici Adi} & \textbf{Sifre} & \textbf{Rol} & \textbf{Fakulte} \\
\hline
superadmin & super123 & Super Admin & - \\
\hline
muh\_admin & admin123 & Fakulte Admin & MUH \\
\hline
muh\_sup1 & super123 & Supervisor & MUH \\
\hline
muh\_temiz1 & temiz123 & Staff & MUH \\
\hline
\end{tabular}
\end{table}

\section{API Test Ornekleri}
\label{sec:api-test}

\begin{lstlisting}[basicstyle=\ttfamily\small, frame=single, 
                   caption={Giris API Istegi}]
POST /api/auth/login
Content-Type: application/json

{
    "username": "superadmin",
    "password": "super123"
}

// Basarili yanit:
{
    "message": "Giris basarili",
    "user": {
        "id": 1,
        "username": "superadmin",
        "role": "super_admin"
    }
}
\end{lstlisting}
